% ======================
% Glossary
% ======================
\chapter*{Bảng thuật ngữ và từ viết tắt}
\addcontentsline{toc}{chapter}{Bảng thuật ngữ và từ viết tắt}

\begingroup
\small
\renewcommand{\arraystretch}{1.3}
\setlength{\tabcolsep}{8pt}

\begin{longtable}{|p{0.25\textwidth}|p{0.35\textwidth}|p{0.3\textwidth}|}
\hline
\textbf{Thuật ngữ tiếng Việt} & \textbf{Thuật ngữ tiếng Anh} & \textbf{Giải thích} \\
\hline
\endfirsthead

\hline
\textbf{Thuật ngữ tiếng Việt} & \textbf{Thuật ngữ tiếng Anh} & \textbf{Giải thích} \\
\hline
\endhead

\hline
\multicolumn{3}{r}{\textit{(Tiếp trang sau)}} \\
\endfoot

\hline
\endlastfoot

\abbrtarget{CV}Thị giác máy tính & Computer Vision (CV) & Lĩnh vực nghiên cứu về khả năng của máy tính trong việc hiểu và xử lý hình ảnh \\
\hline

\abbrtarget{EDA}Phân tích dữ liệu khám phá & Exploratory Data Analysis (EDA) & Giai đoạn khám phá dữ liệu bằng thống kê mô tả và trực quan hóa trước khi phân tích chính \\
\hline

\abbrtarget{CI}Khoảng tin cậy & Confidence Interval (CI) & Khoảng giá trị biểu diễn độ bất định của một ước lượng thống kê \\
\hline

\abbrtarget{SAMPL}Chuẩn báo cáo thống kê & Statistical Analyses and Methods in the Published Literature (SAMPL) & Bộ hướng dẫn báo cáo phương pháp và kết quả thống kê đầy đủ \\
\hline

\abbrtarget{STROBE}Hướng dẫn báo cáo nghiên cứu quan sát & Strengthening the Reporting of Observational Studies in Epidemiology (STROBE) & Checklist báo cáo cho các nghiên cứu quan sát \\
\hline

\end{longtable}

\endgroup

\clearpage
